\documentclass[12pt,a4paper]{article}

\usepackage[utf8]{inputenc}
\usepackage[T1]{fontenc}
\usepackage[brazil]{babel}
\usepackage{lmodern}
\usepackage{geometry}
\usepackage{microtype}
\usepackage{array}
\usepackage{tabularx}
\usepackage{enumitem}
\usepackage{xcolor}
\usepackage{fancyhdr}

\geometry{top=1.7cm,bottom=1.7cm,left=1.8cm,right=1.8cm}
\setlength{\parindent}{0pt}
\setlength{\parskip}{0.45em}
\setlist[itemize]{leftmargin=*,itemsep=0.25em,topsep=0.2em}
\definecolor{azul}{HTML}{173B57}
\definecolor{cinza}{HTML}{F2F4F7}

\pagestyle{fancy}
\setlength{\headheight}{15pt}
\fancyhf{}
\lhead{Sequência didática de \textit{machine learning}}
\rhead{Aula 7: comparação entre modelos}
\cfoot{\thepage}
\renewcommand{\headrulewidth}{0.3pt}

\newcommand{\linha}{\rule{\linewidth}{0.35pt}}
\newcommand{\bloco}[1]{%
  \vspace{0.25em}
  \noindent\colorbox{cinza}{\parbox{0.96\linewidth}{#1}}%
  \vspace{0.3em}
}

\begin{document}

\begin{center}
  {\LARGE\bfseries\color{azul} Aula 7: comparação entre modelos}\par
  \vspace{0.25em}
  {\large Finalidades, dados, resultados e limites}\par
  \vspace{0.35em}
  {\small Roteiro de aplicação da sequência didática}\par
\end{center}

\bloco{\textbf{Lugar desta aula na pesquisa.} Este é o sétimo encontro da sequência de oito aulas. A atividade retoma regressão linear, regressão logística, k-vizinhos mais próximos, árvore de decisão e k-means para comparar finalidades e modos de produzir resultados. As produções individuais são formativas e não integram o corpus principal de pesquisa.}

\section*{1. Identificação}

\begin{tabularx}{\linewidth}{@{}>{\bfseries}p{4.1cm}X@{}}
Turma: & \rule{8.5cm}{0.35pt} \\
Data: & \rule{8.5cm}{0.35pt} \\
Duração: & 50 minutos \\
Organização: & Exploração individual, quadro coletivo e ficha de comparação. \\
Recursos: & Plataforma com os cinco módulos, computadores do laboratório ou Chromebooks, quadro e uma ficha por estudante. \\
\end{tabularx}

\section*{2. Objetivos da aula}

Ao final do encontro, espera-se que os estudantes comparem previsão numérica, classificação e agrupamento; relacionem dados e tipo de saída à escolha de um modelo; reconheçam que parâmetros e regras podem alterar resultados; justifiquem uma escolha sem depender apenas do nome do algoritmo.

\textbf{Relação curricular:} EM13CO10, EM13CO11 e EM13CO12.

\textbf{Conceitos em foco:} dado, finalidade, entrada, saída, previsão, classificação, agrupamento, parâmetro, regra e limite.

\section*{3. Preparação antes da aula}

\begin{itemize}
  \item Abrir os cinco módulos da plataforma e conferir o funcionamento local dos cenários.
  \item Preparar no quadro três colunas: prever um número, escolher uma categoria e formar grupos.
  \item Confirmar que cada estudante terá acesso individual à plataforma, em computador do laboratório ou Chromebook.
  \item Imprimir uma ficha de comparação por estudante, sem código de participante e sem recolhimento como instrumento de pesquisa.
  \item Separar duas situações inéditas breves para discussão coletiva, sem repetir os enunciados reservados ao instrumento final.
\end{itemize}

\bloco{\textbf{Cuidado didático.} A finalidade da tarefa vem antes do nome do modelo. A turma deve começar perguntando qual saída se deseja produzir e quais dados estão disponíveis, evitando associações decoradas entre uma palavra e um algoritmo.}

\section*{4. Desenvolvimento da aula}

\renewcommand{\arraystretch}{1.18}
\begin{tabularx}{\linewidth}{|>{\centering\arraybackslash}p{1.65cm}|p{5.0cm}|X|}
\hline
\textbf{Tempo} & \textbf{Ação com a turma} & \textbf{Mediação docente} \\
\hline
0--8 min & Retomar os cinco modelos em um quadro comparativo. Nomear a finalidade principal de cada atividade anterior. & Organizar a conversa a partir de três saídas: número previsto, categoria e grupo. Evitar revisão baseada apenas em definições decoradas. \\
\hline
8--18 min & Apresentar situações inéditas breves para cada estudante. Pedir que indiquem qual saída desejam em cada situação. & Perguntar primeiro ``qual resultado queremos?'' e depois ``quais dados podem ajudar?''. Só então discutir um modelo possível. \\
\hline
18--33 min & Confirmar duas situações na plataforma, em projeção ou exploração livre orientada. Retomar controles como inclinação, limiar, valor de $k$, profundidade e centróides. & Comparar o papel de cada controle: alguns mudam uma previsão, outros mudam uma regra de decisão ou a organização dos grupos. \\
\hline
33--43 min & Estudantes preenchem a ficha de comparação e explicam uma escolha para a turma. & Pedir justificativas relacionadas a dados e tipo de saída. Corrigir associações como ``dois grupos sempre exigem regressão logística''. \\
\hline
43--50 min & Síntese oral e preparação para o encontro final. & Informar que o instrumento final pedirá explicações sobre situações novas, sem exigir memorização de definições. \\
\hline
\end{tabularx}

\section*{5. Perguntas para conduzir a conversa}

\begin{itemize}
  \item ``Qual saída essa situação pede?''
  \item ``Quais dados seriam usados?''
  \item ``Há categorias já conhecidas?''
  \item ``Qual regra ou parâmetro pode mudar o resultado?''
  \item ``Que limite aparece nessa escolha?''
\end{itemize}

\section*{6. Evidência a observar}

Observar se os estudantes começam pela finalidade da tarefa, relacionam dados e saída à escolha de um modelo e reconhecem que diferentes modelos podem lidar com classificações de formas distintas. No registro docente, anotar confusões recorrentes, justificativas produtivas e condições de infraestrutura que possam ajudar a interpretar a aula.

\newpage

\begin{center}
  {\LARGE\bfseries\color{azul} Ficha de comparação individual}\par
  \vspace{0.2em}
  {\large Aula 7: escolher pela finalidade da tarefa}\par
  \vspace{0.35em}
  {\small Material formativo para discussão durante a aula}\par
\end{center}

\bloco{\textbf{Orientação.} Preencha a ficha individualmente. Ela serve para organizar a discussão e não será usada como instrumento de pesquisa. Em cada situação, indique primeiro a saída desejada e depois o modelo que parece mais adequado.}

\textbf{Nome do estudante:} \rule{8.5cm}{0.4pt}

\renewcommand{\arraystretch}{1.4}
\begin{tabularx}{\linewidth}{|p{3.7cm}|p{2.3cm}|X|}
\hline
\textbf{Modelo} & \textbf{Saída principal} & \textbf{Como produz ou organiza o resultado} \\
\hline
Regressão linear & & \\
\hline
Regressão logística & & \\
\hline
k-vizinhos mais próximos & & \\
\hline
Árvore de decisão & & \\
\hline
k-means & & \\
\hline
\end{tabularx}

\section*{Situação A}

Uma escola quer estimar quanto de energia uma sala pode consumir amanhã a partir de temperatura e quantidade de pessoas. Qual saída é desejada? Qual modelo parece adequado? Justifique.

\vspace{0.25em}
\linha\par\vspace{0.75em}
\linha

\section*{Situação B}

Uma equipe possui pontos de coleta em um mapa, mas ainda não definiu categorias para esses locais. Qual saída é desejada? Qual modelo parece adequado? Justifique.

\vspace{0.25em}
\linha\par\vspace{0.75em}
\linha

\end{document}
